%=============================================================================
% MÓDULO 04: REVISÃO DE LITERATURA (Expansão Adicional)
%=============================================================================
% Frontiers in Public Health — Revisão temática narrativa-crítica
%   Bloco 1: Viés algorítmico em saúde
%   Bloco 2: Disparidades de diabetes
%   Bloco 3: Métricas de justiça
%   Bloco 4: Estratégias de mitigação
%   Bloco 5: Interseccionalidade
%=============================================================================

\section{Revisão de Literatura}

\subsection{Viés Algorítmico em Saúde}

O viés algorítmico foi documentado em numerosas aplicações de saúde, desde 
predição de risco até imagem diagnóstica \cite{rajkomar2018}. Um estudo 
emblemático de Obermeyer et al. \cite{obermeyer2019} revelou que um 
algoritmo de saúde amplamente utilizado subestimou sistematicamente as 
necessidades de saúde de pacientes negros, afetando milhões de indivíduos 
anualmente. Esse achado destacou como decisões algorítmicas aparentemente 
objetivas podem incorporar e ampliar desigualdades históricas, afetando 
milhões de pacientes anualmente.

O estudo de Obermeyer demonstrou que, para um determinado escore de risco, 
os pacientes negros estavam consideravelmente mais doentes do que os 
pacientes brancos, conforme evidenciado por condições como diabetes, 
hipertensão e doença renal crônica. O algoritmo usava custos de saúde como 
proxy para necessidades de saúde, mas como os pacientes negros historicamente 
tinham menos acesso à saúde e, portanto, custos menores, o modelo subestimou 
sistematicamente suas necessidades de saúde. Esse caso ilustra como 
desigualdades históricas podem ser perpetuadas e amplificadas através de 
decisões algorítmicas.

No diabetes especificamente, Chen et al. \cite{chen2019} demonstraram que 
modelos de predição treinados com prontuários eletrônicos mostraram desempenho 
diferenciado entre grupos raciais, com taxas de falso negativo maiores para 
populações minoritárias. O estudo descobriu que modelos treinados 
predominantemente com dados de pacientes brancos mostraram precisão reduzida 
quando aplicados a pacientes negros e hispânicos, sugerindo que a composição 
dos dados de treinamento influencia significativamente a justiça do modelo.

Da mesma forma, uma revisão sistemática de Rajkomar et al. \cite{rajkomar2018} 
identificou disparidades de desempenho em modelos clínicos de AM entre 
subgrupos demográficos, particularmente em configurações com representação 
desbalanceada. A revisão destacou que muitos estudos clínicos de AM falham 
em relatar métricas de desempenho estratificadas por grupos demográficos, 
tornando difícil avaliar a justiça na literatura existente.

As fontes de viés algorítmico são multifacetadas e interconectadas 
\cite{mehrabi2021}:

\begin{itemize}[leftmargin=*]
    \item \textbf{Viés histórico}: Discriminação passada incorporada nos dados 
          de treinamento, refletindo acesso desigual à saúde e padrões de 
          tratamento. Isso inclui protocolos de tratamento diferenciados, 
          níveis variáveis de cuidado e barreiras sistêmicas que historicamente 
          marginalizaram certas populações.
    \item \textbf{Viés de representação}: Sub-representação de certos grupos 
          em conjuntos de dados de treinamento, levando a generalização 
          pobre para essas populações. Quando os modelos são treinados com 
          dados que não representam adequadamente a diversidade da população 
          alvo, eles podem falhar em capturar as relações complexas que 
          existem em grupos sub-representados.
    \item \textbf{Viés de medição}: Protocolos inconsistentes de coleta de 
          dados entre populações, introduzindo erros sistemáticos. Isso pode 
          incluir diferenças em ferramentas de medição, protocolos e padrões 
          que variam entre configurações clínicas e populações.
    \item \textbf{Viés de agregação}: Modelos padronizados que ignoram 
          diferenças populacionais significativas na apresentação e progressão 
          da doença. Diferentes populações podem apresentar fenótipos 
          distintos da doença, padrões de comorbidade e respostas ao 
          tratamento que não são capturados por modelos agregados.
\end{itemize}

As consequências do viés algorítmico vão além do dano individual ao paciente 
para afetar a saúde pública em escala. Quando algoritmos enviesados são 
implantados em sistemas de apoio à decisão clínica, eles podem 
sistematicamente prejudicar populações inteiras, exacerbando disparidades 
existentes de saúde e minando a confiança nos sistemas de saúde. O efeito 
cumulativo desses vieses pode criar ciclos de retroalimentação onde 
populações desfavorecidas recebem pior cuidado, levando a desfechos mais 
ruins que reforçam ainda mais os vieses nos dados.

\subsection{Disparidades de Diabetes}

O diabetes afeta desproporcionalmente comunidades marginalizadas devido a 
determinantes sociais da saúde \cite{narayan2017}. A Federação Internacional 
de Diabetes \cite{idf2021} relata que países de baixa e média renda suportam 
a maior carga, com taxas de prevalência excedendo as de países de alta 
renda em 2--3 vezes. Dentro de países de alta renda, minorias raciais e 
étnicas consistentemente experimentam prevalência mais alta e piores 
desfechos.

Os Centros de Controle e Prevenção de Doenças \cite{cdc2020} relatam 
disparidades raciais e étnicas significativas nos Estados Unidos, conforme 
resumido na Tabela~\ref{tab:diabetes_disparities}. Essas disparidades 
refletem uma interação complexa de fatores sociais, econômicos e ambientais 
que criam desfechos de saúde desiguais entre populações.

\begin{table}[H]
\centering
\caption{Prevalência de Diabetes por Raça/Etnia nos Estados Unidos 
(Fonte: CDC, 2020)}
\begin{tabular}{lcc}
\toprule
\textbf{Grupo} & \textbf{Prevalência (\%)} & \textbf{Risco Relativo} \\
\midrule
Branco & 7.4 & 1.00 \\
Negro & 11.7 & 1.58 \\
Hispânico & 12.5 & 1.69 \\
Asiático & 8.0 & 1.08 \\
Índio Americano & 14.7 & 1.99 \\
\bottomrule
\end{tabular}
\label{tab:diabetes_disparities}
\end{table}

Essas disparidades são impulsionadas por múltiplos fatores interconectados 
\cite{hill2022}: acesso desigual a serviços de saúde, taxas mais altas 
de obesidade e inatividade física, insegurança alimentar e acesso limitado 
a alimentos saudáveis, estresse crônico de discriminação e status 
socioeconômico mais baixo. Quando os modelos de AM são treinados com dados 
que refletem essas disparidades sem correção apropriada, eles podem aprender 
e perpetuar essas desigualdades, criando um ciclo de retroalimentação que 
marginaliza ainda mais populações já desfavorecidas.

Os determinantes sociais da saúde que impulsionam as disparidades de diabetes 
estão profundamente enraizados em desigualdades estruturais, incluindo 
segregação residencial, disparidades educacionais, discriminação no emprego 
e trauma histórico. Esses fatores criam desvantagem cumulativa que se 
manifesta em maior prevalência de diabetes, diagnóstico mais tardio e 
piores desfechos entre populações marginalizadas. Modelos algorítmicos 
treinados com dados moldados por esses fatores estruturais correm o risco 
de codificar e amplificar essas desigualdades.

\subsection{Métricas de Justiça}

Múltiplas métricas de justiça foram propostas para quantificar o viés 
algorítmico \cite{corbett2018}. As métricas mais comumente utilizadas 
incluem:

\subsubsection{Odds Equalizados}

Os Odds Equalizados, introduzidos por Hardt et al. \cite{hardt2016}, 
exigem que a taxa de verdadeiros positivos (sensibilidade) e a taxa de 
falsos positivos (1 - especificidade) sejam iguais entre os grupos:

\begin{equation}
P(\hat{Y}=1 \mid Y=1, A=a) = P(\hat{Y}=1 \mid Y=1, A=b) \quad \forall a,b
\label{eq:eo_tpr}
\end{equation}

\begin{equation}
P(\hat{Y}=1 \mid Y=0, A=a) = P(\hat{Y}=1 \mid Y=0, A=b) \quad \forall a,b
\label{eq:eo_fpr}
\end{equation}

\noindent onde $\hat{Y}$ é o desfecho previsto, $Y$ é o desfecho real 
e $A$ é o atributo sensível. Os Odds Equalizados são particularmente 
relevantes em contextos clínicos onde tanto diagnósticos perdidos (falsos 
negativos) quanto intervenções desnecessárias (falsos positivos) carregam 
consequências significativas.

O critério de Odds Equalizados é especialmente importante na triagem de 
diabetes porque ambos os tipos de erro têm implicações clínicas sérias. 
Falsos negativos podem levar a tratamento atrasado e maior risco de 
complicações, enquanto falsos positivos podem resultar em ansiedade 
desnecessária, testes adicionais e potencial sobretratamento. Ao exigir 
taxas de erro iguais entre os grupos, os Odds Equalizados garantem que 
nenhuma população suporte uma carga desproporcional desses erros.

\subsubsection{Paridade Demográfica}

A Paridade Demográfica exige que a predição seja estatisticamente 
independente do atributo sensível \cite{dwork2012}:

\begin{equation}
P(\hat{Y}=1 \mid A=a) = P(\hat{Y}=1 \mid A=b) \quad \forall a,b
\label{eq:dp}
\end{equation}

Embora mais simples de implementar, a Paridade Demográfica não leva em 
conta diferenças legítimas nas taxas básicas entre os grupos, o que pode 
levar a conclusões paradoxais. Por exemplo, se um grupo tem uma prevalência 
genuinamente mais alta de diabetes, exigir taxas de predição iguais pode 
levar a sub-diagnóstico nesse grupo.

A tensão entre Paridade Demográfica e Odds Equalizados destaca a 
complexidade de definir justiça em contextos clínicos. Diferentes critérios 
de justiça podem ser apropriados dependendo da aplicação clínica específica, 
das consequências de diferentes tipos de erros e dos padrões subjacentes 
de prevalência entre as populações.

\subsubsection{Calibração}

Um modelo está bem calibrado se as probabilidades previstas correspondem 
às taxas reais de desfecho entre os grupos \cite{chouldechova2017}:

\begin{equation}
P(Y=1 \mid \hat{P}=p, A=a) = P(Y=1 \mid \hat{P}=p, A=b) \quad \forall a,b,p
\label{eq:cal}
\end{equation}

Chouldechova \cite{chouldechova2017} demonstrou que, exceto em casos 
triviais, é impossível satisfazer simultaneamente calibração, odds 
equalizados e paridade preditiva, necessitando consideração cuidadosa 
de qual critério de justiça é mais apropriado para o contexto clínico.

A calibração é particularmente importante na decisão clínica porque 
garante que as probabilidades previstas sejam significativas e acionáveis. 
Um modelo bem calibrado permite que os clínicos interpretem escores de 
risco consistentemente entre grupos de pacientes, permitindo uma decisão 
clínica mais equitativa. No entanto, alcançar calibração entre os grupos 
pode exigir ajustes específicos do grupo que levantam considerações 
adicionais de justiça.

\subsection{Estratégias de Mitigação de Viés}

As estratégias de mitigação de viés podem ser categorizadas por sua posição 
no pipeline de AM \cite{friedler2019}:

\begin{enumerate}[leftmargin=*]
    \item \textbf{Pré-processamento} (nivel de dados): Modificar os dados 
          de treinamento para remover viés antes do treinamento do modelo. 
          As técnicas incluem reposicionamento (atribuir diferentes pesos 
          às amostras), reamostragem (sobreamostrar grupos sub-representados) 
          e aumento de dados (gerar amostras sintéticas usando SMOTE ou 
          métodos similares). Os métodos de pré-processamento são atraentes 
          porque são agnósticos ao modelo e podem ser aplicados a qualquer 
          conjunto de dados sem modificar o algoritmo de aprendizado.
    
    \item \textbf{Durante o processamento} (nivel do algoritmo): Modificar 
          o algoritmo de aprendizado para incorporar restrições de justiça. 
          As abordagens incluem dessensibilização adversarial (remover 
          informações sensíveis das representações aprendidas), otimização 
          com restrições de justiça e métodos de regularização que penalizam 
          predições injustas. Os métodos durante o processamento oferecem a 
          vantagem de otimizar conjuntamente desempenho e justiça, mas 
          exigem acesso ao algoritmo de aprendizado.
    
    \item \textbf{Pós-processamento} (nivel de predição): Modificar as 
          predições do modelo para alcançar justiça. Os métodos incluem 
          ajuste de limiar (aplicar diferentes limiares de decisão por grupo), 
          calibração (equalizar probabilidades previstas) e opção de rejeição 
          (abster-se de predição em casos incertos). Os métodos de 
          pós-processamento são flexíveis e podem ser aplicados a qualquer 
          modelo treinado, mas podem não abordar as causas raiz do viés.
\end{enumerate}

Cada categoria oferece vantagens e trade-offs distintos. Os métodos de 
pré-processamento são agnósticos ao modelo mas podem reduzir a precisão 
geral. Os métodos durante o processamento otimizam justiça e precisão 
conjuntamente mas exigem modificação algorítmica. Os métodos de 
pós-processamento são flexíveis mas podem não abordar as causas raiz do 
viés. Na prática, uma combinação de abordagens pode ser mais eficaz, 
diferentes fontes de viés podem exigir diferentes estratégias de mitigação.

\subsection{Interseccionalidade na Justiça Algorítmica}

A teoria da interseccionalidade, originalmente proposta por Crenshaw 
\cite{crenshaw1989}, destaca como múltiplas formas de discriminação 
interagem para criar modos únicos de desvantagem. Na justiça algorítmica, 
a análise interseccional examina viés na intersecção de múltiplos 
atributos sensíveis (por exemplo, raça \textit{e} gênero) em vez de cada 
atributo independentemente \cite{buolamwini2018}.

Buolamwini e Gebru \cite{buolamwini2018} demonstraram esse fenômeno em 
sistemas de reconhecimento facial, onde classificadores comerciais 
apresentaram disparidades significativas de desempenho para mulheres de 
pele escura, apesar de precisão razoável para indivíduos de pele clara 
e homens. Essa descoberta destacou as limitações da análise de justiça 
de atributo único e a importância de examinar efeitos interseccionais.

Em contextos de saúde, o viés interseccional é particularmente preocupante 
porque as disparidades de saúde frequentemente se agrupam ao longo de 
múltiplas dimensões. Indivíduos que são simultaneamente membros de múltiplos 
grupos marginalizados podem experimentar desvantagem acumulada que não é 
capturada analisando cada atributo isoladamente. Por exemplo, mulheres 
negras podem enfrentar discriminação racial e de gênero na saúde, criando 
barreiras únicas que não são capturadas analisando raça ou gênero 
separadamente.

Apesar de sua importância, a análise interseccional continua sub-representada 
na pesquisa sobre justiça de IA em saúde, representando uma lacuna crítica 
que este estudo aborda. Os poucos estudos existentes que examinaram efeitos 
interseccionais em IA em saúde encontraram disparidades significativas na 
intersecção de raça e gênero, sugerindo que análises de atributo único 
podem subestimar sistematicamente a extensão do viés algorítmico.

O arcabouço da interseccionalidade também destaca a importância de considerar 
fatores estruturais e contextuais que moldam as disparidades de saúde. O 
viés algorítmico não ocorre isoladamente, mas está embutido em sistemas mais 
amplos de desigualdade que influenciam a coleta de dados, o desenvolvimento 
de modelos e a implantação. Compreender esses fatores estruturais é essencial 
para desenvolver estratégias eficazes para abordar o viés algorítmico em saúde.
